\documentclass[11pt]{amsart}
\usepackage[T1]{fontenc}
\usepackage{lmodern}
\usepackage{amsmath,amssymb,amsthm,mathtools}
\usepackage{booktabs}
\usepackage{microtype}
\usepackage{xcolor}
\usepackage[colorlinks=true,linkcolor=blue!55!black,citecolor=blue!55!black,urlcolor=blue!55!black]{hyperref}

\newtheorem{theorem}{Theorem}[section]
\newtheorem{lemma}[theorem]{Lemma}
\newtheorem{proposition}[theorem]{Proposition}
\theoremstyle{definition}
\newtheorem{definition}[theorem]{Definition}
\newtheorem{remark}[theorem]{Remark}

\newcommand{\Lyn}{\mathcal{L}}
\newcommand{\ov}{\operatorname{ov}}
\newcommand{\pre}{\operatorname{pre}}
\newcommand{\suf}{\operatorname{suf}}

\title[Shortest superstring of length-eight Lyndon words]{The shortest common superstring of the length-eight binary Lyndon words has length 94}
\author{Zijian Zeng}
\address{Institute of Computer Science and Digital Innovation, UCSI University, Kuala Lumpur 56000, Malaysia}
\email{zijianzeng@foxmail.com}
\email{1002266693@ucsiuniversity.edu.my}
\date{August 30, 2026}
\hypersetup{
  pdftitle={The shortest common superstring of the length-eight binary Lyndon words has length 94},
  pdfauthor={Zijian Zeng}
}

\begin{document}

\begin{abstract}
There are 30 binary Lyndon words of length eight.  We determine the minimum
length of a binary word containing all of them as contiguous factors.  The
answer is 94.  The lower bound orders selected target occurrences and relaxes
the resulting maximum-overlap Hamilton path to a 31-node assignment problem.
An integer assignment dual of value 146 proves the lower bound
$30\cdot8-146=94$, while an explicit 94-bit word attains it.  The finite dual,
all target occurrences, and two independent exact verification procedures are
provided.
\end{abstract}

\maketitle

\section{Introduction}

A Lyndon word is the lexicographically least word in a primitive cyclic
rotation class.  We consider the following finite shortest-common-superstring
problem: determine the shortest binary word containing every binary Lyndon
word of length eight as a contiguous factor.  In the notation below, its
minimum length is $L_8$.

The reviewed TheoremDB packet for this problem recorded the rigorous interval
$49\le L_8\le94$: a direct Boolean satisfiability computation excluded length
48, and a replayable 94-bit word supplied the upper bound
\cite{TheoremDBP2600}.  That packet did not retain a proof-producing SAT trace
and listed the exact value as open.  The classical Lyndon-concatenation
construction produces a de Bruijn sequence for a larger target family
\cite{FredricksenMaiorana1978}; it does not settle this restricted instance.
The standard reduction from shortest superstrings to maximum-overlap
asymmetric paths is described, for example, by Mucha
\cite{Mucha2013}.

We prove that the recorded upper endpoint is optimal.  The main point is that
the assignment relaxation of the 30-target maximum-overlap path is already
tight.  Its optimum is bounded above by a short integer dual certificate of
value 146, and a target order attains that value.  The proof therefore needs
neither a floating-point optimizer nor a solver's uncheckable status report.

Our exact-title, endpoint, and full-witness searches found no public exact
solution in the sources searched through August 30, 2026.  This is a
search-bounded statement, not a claim about inaccessible or unpublished work.

\section{Definitions and the occurrence-order bound}

Let $\Sigma=\{0,1\}$ with $0<1$.  For
$x=x_0x_1\cdots x_7\in\Sigma^8$, put
\[
 \rho_j(x)=x_jx_{j+1}\cdots x_7x_0\cdots x_{j-1}
 \qquad(1\le j\le7).
\]
The target set is
\[
 \Lyn_8=\{x\in\Sigma^8:x<\rho_j(x)\text{ for every }1\le j\le7\}.
\]
All superstrings in this paper are finite linear words; there is no cyclic
wraparound.

The number of primitive binary words of length eight is, by M\"obius
inversion,
\[
 \sum_{d\mid8}\mu(d)2^{8/d}=2^8-2^4=240.
\]
Each primitive rotation orbit has size eight and a unique lexicographically
least representative.  Thus $|\Lyn_8|=240/8=30$.

For distinct $x,y\in\Sigma^8$, define their directed overlap by
\[
 \ov(x,y)=\max\{k\in\{0,\ldots,7\}:\suf_k(x)=\pre_k(y)\}.
\]

\begin{lemma}[Occurrence-order bound]\label{lem:occurrence}
For every word $W$ containing every member of $\Lyn_8$, there is an ordering
$x_1,\ldots,x_{30}$ of $\Lyn_8$ such that
\[
 |W|\ge 240-\sum_{r=1}^{29}\ov(x_r,x_{r+1}).
\]
\end{lemma}

\begin{proof}
Choose one occurrence of each target and sort their starting positions:
$p_1<\cdots<p_{30}$.  Distinct length-eight factors cannot have the same
starting position, so this gives an ordering $x_1,\ldots,x_{30}$.
If $p_{r+1}-p_r<8$, the two occurrences force a suffix-prefix match of length
$8-(p_{r+1}-p_r)$; if the displacement is at least eight, the same inequality
is immediate.  Hence
\[
 p_{r+1}-p_r\ge8-\ov(x_r,x_{r+1}).
\]
Finally, $|W|\ge p_{30}+8\ge p_{30}-p_1+8$.  Summing the 29 displacement
inequalities proves the claim.
\end{proof}

\section{The exact assignment certificate}

Adjoin a dummy node $*$ to $\Lyn_8$.  For distinct nodes $i,j$, define
\[
 c_{ij}=\begin{cases}
 \ov(i,j),&i,j\in\Lyn_8,\\
 0,&i= *\text{ or }j= *.
 \end{cases}
\]
An ordered target path $x_1,\ldots,x_{30}$, together with the arcs
$*\to x_1$ and $x_{30}\to *$, is a permutation assignment on the 31 nodes.
Consequently, any vectors $(u_i)$ and $(v_j)$ satisfying
\begin{equation}\label{eq:dual}
 u_i+v_j\ge c_{ij}\qquad(i\ne j)
\end{equation}
bound every target-path overlap sum by $\sum_i u_i+\sum_jv_j$.  This statement
does not require the assignment itself to be connected; allowing disconnected
assignments only relaxes the maximization problem.

Table~\ref{tab:certificate} lists a feasible integer dual.  It also lists a
zero-based occurrence position $p$ for every target in the word used for the
upper bound.  The dummy has $(u_*,v_*)=(0,0)$.  Direct suffix-prefix comparison
checks all $31\cdot30=930$ inequalities in \eqref{eq:dual}.  The row sums are
\[
 \sum_i u_i=45,\qquad \sum_jv_j=101.
\]

\begin{table}[ht]
\centering
\caption{Integer assignment-dual values and target occurrence positions.
The two halves continue one lexicographically ordered list.}
\label{tab:certificate}
\scriptsize
\setlength{\tabcolsep}{3.4pt}
\begin{tabular}{@{}lrrr@{\qquad}lrrr@{}}
\toprule
target & $u$ & $v$ & $p$ & target & $u$ & $v$ & $p$\\
\midrule
\texttt{00000001}&6&0&80 & \texttt{00100101}&0&0&68 \\
\texttt{00000011}&6&1&81 & \texttt{00100111}&0&3&15 \\
\texttt{00000101}&6&0&45 & \texttt{00101011}&0&7&36 \\
\texttt{00000111}&6&1&82 & \texttt{00101101}&0&5&71 \\
\texttt{00001001}&6&0&13 & \texttt{00101111}&0&7&48 \\
\texttt{00001011}&1&1&46 & \texttt{00110101}&0&0&5 \\
\texttt{00001101}&6&0&57 & \texttt{00110111}&0&7&59 \\
\texttt{00001111}&1&1&83 & \texttt{00111011}&0&7&27 \\
\texttt{00010011}&4&1&14 & \texttt{00111101}&0&5&18 \\
\texttt{00010101}&0&0&35 & \texttt{00111111}&0&7&85 \\
\texttt{00010111}&0&6&47 & \texttt{01010111}&0&7&37 \\
\texttt{00011001}&3&0&0 & \texttt{01011011}&0&7&72 \\
\texttt{00011011}&0&1&58 & \texttt{01011111}&0&7&49 \\
\texttt{00011101}&0&0&26 & \texttt{01101111}&0&7&60 \\
\texttt{00011111}&0&6&84 & \texttt{01111111}&0&7&86 \\
\bottomrule
\end{tabular}
\end{table}

For completeness, one target order attaining overlap 146 is
\begin{align*}
&\texttt{00011001},\texttt{00110101},\texttt{00001001},
\texttt{00010011},\texttt{00100111},\texttt{00111101},\\
&\texttt{00011101},\texttt{00111011},\texttt{00010101},
\texttt{00101011},\texttt{01010111},\texttt{00000101},\\
&\texttt{00001011},\texttt{00010111},\texttt{00101111},
\texttt{01011111},\texttt{00001101},\texttt{00011011},\\
&\texttt{00110111},\texttt{01101111},\texttt{00100101},
\texttt{00101101},\texttt{01011011},\texttt{00000001},\\
&\texttt{00000011},\texttt{00000111},\texttt{00001111},
\texttt{00011111},\texttt{00111111},\texttt{01111111}.
\end{align*}
Appending each target after deleting its maximal overlap with its predecessor
gives the witness below.

\section{Exact value}

\begin{theorem}\label{thm:main}
The minimum length of a binary word containing all 30 binary Lyndon words of
length eight as contiguous factors is
\[
 \boxed{L_8=94}.
\]
\end{theorem}

\begin{proof}
By Lemma~\ref{lem:occurrence}, every common superstring has length at least
$240$ minus the overlap weight of some target path.  Completing that path
with the dummy gives an assignment.  Summing \eqref{eq:dual} over its 31 arcs
and using Table~\ref{tab:certificate} bounds the overlap by $45+101=146$.
Therefore every common superstring has length at least
\[
 240-146=94.
\]

For the reverse inequality, concatenate the following two displayed lines
without a separator:
\begin{center}
\small\ttfamily
00011001101010000100111101000111011000101011100\\
00010111110000110111100100101101100000001111111
\end{center}
Each line has 47 bits, so the resulting word has length 94.  The positions in
Table~\ref{tab:certificate} verify that every target occurs.  Hence $L_8\le94$,
which matches the lower bound.
\end{proof}

\section{Exact verification and limitations}

The supplementary certificate stores the witness, the attaining target order,
and the integer dual.  A standard-library Python verifier regenerates the
targets from strict rotation comparisons, checks all occurrences, rebuilds
the witness from the target order, and checks all 930 dual inequalities.
A second standard-library program uses a different target generator based on
cyclic orbits, verifies coverage by rolling windows, ignores the stored dual
and order, and recomputes the assignment optimum with an integer Hungarian
algorithm.  The latter implementation was also compared with exhaustive
permutation search on deterministic small test matrices.  No proof assistant
was used.

The result is confined to the specified finite parameter $n=8$.  It does not
give a formula for other word lengths, nor does it establish that the
assignment relaxation remains tight for another target family.

\section*{Data, code, and exact verification}

The accompanying source archive includes the JSON certificate, both exact
checkers, regression tests, and a diagnostic discovery script.  From the
archive root the exact checks are run as follows:
\begin{center}
\begin{minipage}{0.92\linewidth}
\ttfamily\small
python3 verify/verify\_certificate.py\\
python3 verify/independent\_checker.py\\
python3 verify/test\_verifiers.py
\end{minipage}
\end{center}
The proof depends only on integer and string operations in the exact checkers.
A mixed-integer optimizer was used to locate a candidate path, but no
floating-point objective, optimizer status, or timeout is used in the proof.

\bibliographystyle{amsplain}
\bibliography{references}

\end{document}
